\documentclass[11pt]{article}
\usepackage[utf8]{inputenc}
\usepackage[T1]{fontenc}
\usepackage{lmodern}
\usepackage[margin=1in]{geometry}
\usepackage{amsmath,amssymb,amsthm}
\usepackage{enumitem}
\usepackage{booktabs}
\usepackage{xcolor}
\usepackage{tikz}
\usepackage[colorlinks=true,bookmarks=false,linkcolor=blue!50!black,urlcolor=blue!50!black]{hyperref}

\theoremstyle{plain}
\newtheorem{theorem}{Theorem}
\newtheorem{lemma}[theorem]{Lemma}
\newtheorem{claim}[theorem]{Claim}
\newtheorem{corollary}[theorem]{Corollary}
\theoremstyle{definition}
\newtheorem{definition}[theorem]{Definition}
\newtheorem{example}[theorem]{Example}
\theoremstyle{remark}
\newtheorem*{remark}{Remark}

\newcommand{\N}{\mathbb{N}}
\newcommand{\Z}{\mathbb{Z}}
\newcommand{\R}{\mathbb{R}}
\newcommand{\Q}{\mathbb{Q}}
\newcommand{\iv}[1]{[#1]}                 % Iverson bracket
\newcommand{\powerset}{\mathcal{P}}
\DeclareMathOperator{\ord}{ord}
\DeclareMathOperator{\gcdop}{gcd}

\setlength{\parindent}{0pt}
\setlength{\parskip}{0.6em}
\newcommand{\lesson}[2]{\href{../lessons/#1}{#2}}   % link to a lesson image

% Let long inline lists like \{1,2,5,6\} break across lines after a comma.
\mathchardef\mycomma="613B
\mathcode`,="8000
{\catcode`,=\active \gdef,{\mycomma\penalty0\hspace{0pt}}}
\emergencystretch=2em

\title{Counting and double counting}
\author{Notes from Carlos' lessons}
\date{}
\begin{document}
\maketitle

\textbf{Source lessons:} \lesson{2026-02-11-proving-bijections.png}{2026-02-11 (proving bijections)};
pigeonhole fragment at the left of \lesson{2026-02-27-ramsey-numbers.png}{2026-02-27 (Ramsey numbers)}.

\section*{Iverson bracket}

For a proposition $P$, $\iv{P}=1$ if $P$ is true and $0$ otherwise. Rules (check by truth table):
\begin{enumerate}[nosep]
  \item $\iv{P\wedge Q}=\iv{P}\iv{Q}$
  \item $\iv{P\vee Q}=\iv{P}+\iv{Q}-\iv{P}\iv{Q}$
  \item $\iv{\neg P}=1-\iv{P}$
\end{enumerate}

\emph{Cardinality as a sum:} for finite $X$ and $S\subseteq X$, $|S|=\sum_{x\in X}\iv{x\in S}$.
Hence $n-|S|=\sum_{t\in[n]}\bigl(1-\iv{t\in S}\bigr)=\sum_{t\in[n]}\iv{t\notin S}$ (used below).

The technique of the whole lesson: write a count as a double sum of brackets, swap the order of summation, and evaluate the other way.

\section*{Handshake lemma}

\begin{theorem}[Handshake lemma]
For a finite undirected graph $G=(V,E)$ (edges are $2$-element subsets of $V$),
$\displaystyle\sum_{v\in V}\deg(v)=2|E|$.
\end{theorem}
\begin{proof}
$\deg(v)=|\{e\in E: v\in e\}|=\sum_{e\in E}\iv{v\in e}$. So
\[
\sum_{v}\deg(v)=\sum_{v}\sum_{e}\iv{v\in e}=\sum_{e}\sum_{v}\iv{v\in e}
=\sum_{e}\bigl|\{v: v\in e\}\bigr|=\sum_{e}2=2|E|.
\]
The swap is legal because both sums are finite. The set being double counted is $\{(v,e)\in V\times E: v\in e\}$.
\end{proof}

\begin{corollary}
The number of odd-degree vertices is even; in particular a graph on an odd number of vertices cannot have all degrees odd.
\end{corollary}
(This is used in the parity argument for $R(3,4)\le 9$.)

\section*{$\sum_k\binom nk = 2^n$}

\begin{theorem}
$\displaystyle\sum_{k=0}^n\binom nk=2^n$.
\end{theorem}
\begin{proof}
$\binom nk=|\{S\subseteq[n]:|S|=k\}|=\sum_{S\subseteq[n]}\iv{|S|=k}$. Sum over $k=0,\dots,n$ and swap:
\[
\sum_{k=0}^n\binom nk=\sum_{S\subseteq[n]}\sum_{k=0}^n\iv{|S|=k}=\sum_{S\subseteq[n]}1=|\powerset([n])|=2^n,
\]
since for each $S$ exactly one $k$ (namely $|S|$) makes the bracket equal to $1$.
\end{proof}

\section*{Binomial theorem}

\begin{theorem}
$(x+y)^n=\displaystyle\sum_{k=0}^n\binom nk x^ky^{n-k}$.
\end{theorem}
\begin{proof}
Expand the product $\prod_{t=1}^n(x+y)$. Choosing $x$ or $y$ from each factor $t$ corresponds to choosing the set $S\subseteq[n]$ of factors that contribute $x$. So
\[
(x+y)^n=\sum_{S\subseteq[n]}\ \prod_{t\in[n]}x^{\iv{t\in S}}y^{\iv{t\notin S}}
=\sum_{S\subseteq[n]}x^{|S|}y^{\,n-|S|}.
\]
Group the sets $S$ by size $k=|S|$; there are $\binom nk$ of each size.
\end{proof}

The board example: $n=3$, $S=\{1,3\}$ picks $x\cdot y\cdot x$.
A second proof, via exponential generating functions, is in the topic \emph{Exponential generating functions}.

\section*{Pigeonhole: a block of terms divisible by $N$}

The left margin of the 02-27 board asks: given $a_1,\dots,a_N\in\N$, is there a nonempty $S$ with $N\mid\sum_{k\in S}a_k$? It lists the prefix sums $a_1,\ a_1+a_2,\ \dots,\ a_1+\dots+a_N$, which is the standard proof.

\begin{claim}
Some nonempty block of \emph{consecutive} terms has sum divisible by $N$.
\end{claim}
\begin{proof}
Let $s_j\equiv a_1+\dots+a_j\pmod N$ for $j=1,\dots,N$. If some $s_j=0$, the block $a_1,\dots,a_j$ works. Otherwise the $N$ values lie in $\{1,\dots,N-1\}$, so two coincide, $s_i=s_j$ with $i<j$, and then $a_{i+1}+\dots+a_j\equiv0\pmod N$.
\end{proof}

\section*{Notes}
\begin{itemize}[nosep]
  \item The lesson's filename says ``bijections'', but what is on the board is the Iverson-bracket / double-counting toolkit; no explicit bijection is proved there.
  \item The right-hand scrawl on the board (the $\prod x^{\iv{t\in S}}y^{\iv{t\notin S}}$ expansion and the $S=\{1,3\}$ example) is the binomial-theorem proof above.
\end{itemize}

\end{document}
